\documentclass[10pt]{article}
\usepackage[margin=20mm]{geometry}
\usepackage{amsmath,amssymb,mathtools,fontspec,unicode-math,xurl,hyperref}
\IfFontExistsTF{DejaVu Serif}{\setmainfont{DejaVu Serif}}{\setmainfont{Latin Modern Roman}}
\setmathfont{latinmodern-math.otf}
\setlength{\parindent}{0pt}\setlength{\parskip}{6pt}\sloppy
\title{A verification of selected arguments in
Homomorphisms to 3–Manifold Groups}
\author{}
\date{}
\begin{document}
\maketitle
Abstract: The reductions to equational noetherianity, the stabilization argument for epimorphisms, and selected finite-generation and resolution arguments are examined in the supplied author manuscript. One local auxiliary embedding claim is false: a permitted finite cyclic group need not embed in \(\mathrm{SL}(3,\mathbb Z)\). An explicit uniform embedding into \(\mathrm{GL}(3,\mathbb C)\) repairs that use of linearity without changing the subsequent group-theoretic conclusions. No further substantive defect is established within the declared scope. This is an initial technical review; scientific adjudication and independent acceptance remain pending.\par
\section*{Identity and scope}
The reviewed work is Daniel Groves, Michael Hull, and Hao Liang, Homomorphisms to 3–Manifold Groups, arXiv:2103.06095v3, 12 July 2024 (the title-page date is 15 July 2024), 55 physical PDF pages, including Appendices A and B [paper]. The assigned journal identifier is DOI 10.1007/s00039-025-00713-1. Page references use one-based physical PDF pages, which agree with the printed numbering. The acquired preprint is not identified with an unexamined version of record.\par
The review checks the logical reductions on pp. 1–15, especially Theorems 2.4–2.8, Proposition 3.4 and the first part of Theorem 3.23; it also follows selected non-divergence, resolution, and edge-twisted arguments in Sections 6 and 9 and Appendix A. It does not claim an exhaustive reproof of the collapsed-space geometry, the shortening theorem, the relative JSJ theory, or the external results cited by the authors. The detailed ledger distinguishes those dependencies from arguments checked here. No author code was required or executed.\par
\section*{Definitions and checked reductions}
An equationally noetherian family \(\mathcal G\) has the following uniform property: for every finitely generated \(G\), a finitely presented \(G_0\) admits an epimorphism \(\alpha:G_0\to G\) such that \(\alpha^*:\operatorname{Hom}(G,\Gamma)\to
\operatorname{Hom}(G_0,\Gamma)\) is bijective for every \(\Gamma\in\mathcal G\). With the fixed nonprincipal ultrafilter \(\omega\), the stable kernel is \[K_\omega=\{g:\phi_i(g)=1\ \text{for $\omega$-almost every }i\},
 \qquad L=G/K_\omega .\] Theorem 2.4 equates the family property with \(\omega\)-almost-sure factorization through this quotient. Its external proof is a declared input, rather than a result reproduced here.\par
\subsection*{The descending-chain consequence}
For the surjections in Theorem B, put \(\rho_i=\tau_i\cdots\tau_1\) and \(K_i=\ker\rho_i\). The composition law gives \(K_i\subseteq K_{i+1}\) and \(K_\omega=\bigcup_i K_i\). Factorization of one \(\rho_j\) through the limit gives \(K_\omega\subseteq K_j\), hence equality and stabilization. Subsequent surjections are injective and therefore isomorphisms. This verifies the argument on pp. 5–6. The reversed inclusion printed on p. 5 is resolved by the definitions and the immediately following ascending-kernel argument; it is not counted as a separate substantive finding.\par
\subsection*{Uniformity in finite-index passage}
In Proposition 3.4 (p. 7), a finitely generated \(G\) has finitely many subgroups of index at most a fixed \(n\). Their intersection \(H\) has finite index and is finitely generated. Each \(\phi_i(H)\) lies in the selected subgroup of the target. Factorization on \(H\), followed by a finite intersection of \(\omega\)-large sets for coset representatives, gives factorization on \(G\). The finite bound on the index is essential: the argument does not assert a uniform result over unbounded indices.\par
For the central extension in Corollary 3.18 (p. 9), finite presentation of the quotient and finite generation of the total group give finitely many normal generators for the central kernel; centrality turns these into ordinary generators. The resulting finite generation and finite-presentation deductions are consistent with the stated hypotheses. Lemma 3.19 also retains its centrality requirement: the preimage of a cyclic group in this central extension is abelian.\par
\section*{Problem 1: an integral linear target is too restrictive}
Assessment: Minor, auxiliary proof statement. The first paragraph of the proof of Theorem 3.23 on p. 10 reduces five geometries to subgroups of index at most \(240\) that are cyclic, \(\mathbb Z^3\), or \(\mathbb Z^2\rtimes\mathbb Z\), and asserts that each such group embeds in \(\mathrm{SL}(3,\mathbb Z)\). The actual PDF has \(\mathbb Z\), so this is not an extraction error. Lemma 3.1 (p. 6), used in the next step, only needs one fixed \(\mathrm{GL}(n,K)\).\par
\subsection*{An admissible counterexample}
Let \(\zeta=\exp(2\pi i/241)\) and let \(C_{241}\) act on \(S^3\subset\mathbb C^2\) by \[(z_1,z_2)\longmapsto(\zeta z_1,\zeta z_2).\] Every nonidentity power acts freely, since a point of \(S^3\) has a nonzero coordinate. Thus the spherical lens space \(L(241,1)\) is a closed geometric 3-manifold with fundamental group \(C_{241}\). The integer \(241\) is prime: none of \(2,3,5,7,11,13\), the primes not exceeding \(\sqrt{241}\), divides it. Its only proper subgroup has index \(241\). Consequently passage to a subgroup of index at most \(240\) cannot remove this example.\par
If \(C_{241}\) embedded even in \(\mathrm{GL}(3,\mathbb Q)\), the image \(A\) of a generator would have order \(241\). Since \(A^{241}=I\) and \(A\ne I\), some eigenvalue would be a primitive \(241\)st root of unity. Its minimal polynomial over \(\mathbb Q\) is \[\Phi_{241}(x)=1+x+\cdots+x^{240},\] which is irreducible: \(\Phi_{241}(x+1)=((x+1)^{241}-1)/x\) satisfies Eisenstein’s criterion at \(241\). Its degree \(240\) exceeds the degree of the characteristic polynomial of a \(3\times3\) rational matrix. This contradiction excludes an embedding into \(\mathrm{SL}(3,\mathbb Z)\) in particular. It verifies every group and geometry condition used in the stated reduction.\par
\subsection*{A uniform repair and its precise effect}
The required argument works in the single fixed target \(\mathrm{GL}(3,\mathbb C)\). For a finite cyclic group of order \(m\), use \(\operatorname{diag}(\exp(2\pi i/m),1,1)\). For \(\mathbb Z^3\), use \[(a,b,c)\longmapsto
 \operatorname{diag}(2^a3^b5^c,1,1),\] which is injective by unique prime factorization. Infinite cyclic groups are included by the analogous map \(k\mapsto
\operatorname{diag}(2^k,1,1)\).\par
Equation reference: eq:repair. For any \(A\in\mathrm{GL}(2,\mathbb Z)\), write the multiplication in \(\mathbb Z^2\rtimes_A\mathbb Z\) as \((v,k)(w,l)=(v+A^k w,k+l)\). Then \[
 (v,k)\longmapsto
 2^k
 \begin{pmatrix}
 A^k&v\\
 0&1
 \end{pmatrix}
 \quad\text{in }\mathrm{GL}(3,\mathbb C)\] is a homomorphism by block multiplication. Its lower-right entry recovers \(k\), and with \(k=0\) its upper-right column recovers \(v\), so it is injective. The scalar factor matters: an unweighted affine representation alone can lose the \(\mathbb Z\) factor when \(A\) has finite order.\par
Apply Lemma 3.1 with this fixed dimension and field, then the already stated bounded-index result, Proposition 3.4. This restores the first paragraph’s equational-noetherian conclusion under its same classification input. No new uniformity in \(m\), \(A\), or the index is required, and the bound \(240\) is unchanged.\par
The false integral embedding is a genuine local proof assertion. The repair explicitly supplies the linearity actually required, so the finding does not establish failure of Theorem 3.23, Theorem 2.5, or Theorem A. Its role is an auxiliary geometric-family reduction; the new collapsed-space and resolution mechanisms are not invalidated by this example. This limited dependency and the verified replacement, rather than repair length alone, justify the Minor assessment. No assertion is made about whether a later published version already changes the coefficient ring.\par
\section*{Finite generation and resolution: checked scope}
Theorem 2.7 (p. 13) needs finite generation of the abelian edge groups before applying the vertex-group factorization argument; it does not assume the desired conclusion without comment. For Theorem 2.8 (p. 15), the maps in Theorem 4.20 are injective on stably parabolic subgroups. Their composition embeds each such subgroup in a stably parabolic subgroup of the terminal non-\(\mathcal C\)-divergent limit. Theorem 4.19 makes the latter finitely generated, and Lemma 4.6 makes it abelian. A subgroup of a finitely generated abelian group is finitely generated, completing the stated reduction. Without abelianness that last implication would be invalid; the hypothesis is explicitly present here.\par
The selected Section 9 check follows the finite-ball diagonal construction, the restricted descending-chain step, and the factorization contradiction. It does not replace the shortening theorem with a general termination assumption. In Appendix A, the checked amalgam and edge-twist calculations support the specified algebraic bridges; the omitted proof of Lemma A.4 and the invoked external geometric and JSJ results remain declared dependencies. No additional counterexample is established from these passages. Full verification of every shortening and collapsed-space estimate is outside this review.\par
\section*{Conclusion}
Dimension | State\par
Issue assessment | One supported local Minor finding\par
Repair verification | Fixed \(\mathrm{GL}(3,\mathbb C)\) construction proved above\par
Main theorems | Not refuted by the finding\par
Author adoption | Not established\par
Scientific decision | Pending; no human review entered\par
Independent acceptance | Pending\par
The source-version limit and the stated unexamined dependencies remain part of this assessment. Optional code was not identified in the bounded acquisition record; its absence is not a mathematical defect.\par
\section*{Scope and references}
The current proof in Problem 1 is an analytic group-theoretic argument, not a finite experiment. No symbolic or numerical program is offered as a substitute for it. The existing source extraction is reused by exact source/text hashes and keeps its historical execution origin. Actual source-page views, the current exact-source compiler run, retained dependencies, PDF parsing, and all report-page views are separately bound in the per-paper verification records. These AI checks are initial work, not independent acceptance or a human attestation.\par
1 D. Groves, M. Hull, and H. Liang, Homomorphisms to 3–Manifold Groups, arXiv:2103.06095v3, 12 July 2024. Assigned journal DOI: 10.1007/s00039-025-00713-1.\par
\end{document}
